\resizebox{\linewidth}{!}{%
\begin{tikzpicture}[bi]
\fill[biRH, opacity=0.12] (3.4,-0.35) rectangle (11.6,0.35);
\draw[->] (0,0) -- (12,0) node[right]{$\operatorname{Re}s$};
\draw (1,0.12) -- (1,-0.12) node[below]{$\tfrac12$};
\draw (11,0.12) -- (11,-0.12) node[below]{$1$};
\draw[biRH, thick] (3.4,0.18) -- (3.4,-0.18);
\node[biRH, font=\scriptsize, anchor=south east] at (3.35,0.25)
  {$\sigma_0=\tfrac12+\tfrac{\mu}{2}$};
\node[bi fail pt] at (5.8,0) {};
\node[biFail, font=\scriptsize, anchor=south west] at (5.9,0.25)
  {hypothetical zero $\rho$};
\draw[biBound, thick, decorate, decoration={brace, amplitude=4pt, mirror}]
  (1,-0.6) -- (5.8,-0.6) node[midway, below=5pt, font=\scriptsize]
  {$\mu=\min(\operatorname{Re}\rho-\tfrac12,\tfrac12)$};
\node[font=\scriptsize, biRH, anchor=north] at (8.9,-0.45) {zero-free if the
  bound holds};
\end{tikzpicture}}

\smallskip
\begin{minipage}{0.9\linewidth}\footnotesize
Chebyshev inequality on windows: $G_{\lev}(x)\le
x^{1/2+(e_D+e_L+7)\eps}$, with $e_D=2$ and $e_L=0$ for the families of
Theorem~\ref{thm:tied-families}, so $x^{1/2+9\eps}$. Choosing
$9\eps\le\mu/2$ makes the exponent at most
$\sigma_0$; the windowed Landau theorem then rules out zeros in the shaded
half-plane $\operatorname{Re}s>\sigma_0$, which contains $\rho$. The closer
$\rho$ is to the line, the smaller $\eps$ must be.
\bischematic
\end{minipage}
